\documentclass[main]{subfiles}
\usepackage[sols]{workbook}
\numbering{ws2-10}

\begin{document}

\ifSubfilesClassLoaded{%
  \chaptermark{\chaptertwo}
}{}

\section{Ratio Test} \label{ws2-10}
\begin{framed}
\centerline{\textbf{Lesson Goals}}

You should be able to:
\begin{itemize}
\item Understand and explain the intuition behind the Ratio Test
\item Apply the Ratio Test to determine if a given series converges or diverges, and determine when the Ratio Test gives no information about series convergence
\item Use the Ratio Test to determine the values of $x$ for which a given Taylor series converges or diverges
\end{itemize}
\end{framed}

\begin{enumerate}
  \item 
  \begin{enumerate}
    \item

    %\note{This was on their homework; it's just here so that we can discuss it in class.  They should see that the number 5 is not mportant and that the statement remains true if 5 is replaced by any number $> 1$.  (They may think that the statement is true even when $5$ is replaced by $1$, which can lead to some interesting discussion.)}
    
    \begin{problem}[1in]
    Suppose you have a series $\displaystyle \sum_{k=1}^\infty a_k$ whose
    terms are nonzero (but the series could have both positive and
    negative terms), and you know that $\displaystyle \lim_{k \to \infty}
    \left| \frac{a_{k+1}}{a_k} \right| = 5$.  What does this tell you
    about the magnitude (absolute value) of the terms as $k$ gets really
    big?  Does the series converge or diverge, or is there not enough
    information to tell?  
    \end{problem}

    \begin{solution}
    For $k$ large enough we have $6>|a_{k+1}/a_k| > 4$. In particular $|a_{k+1}| > 4 |a_k|$. Thus $\lim_{n \to \infty} |a_k| = \infty$ and the series diverges by the $n$th term test for divergence.
    \end{solution}

    \item

    %\note{This was also on their homework and here so that we can discuss it in class.  They should see here that this remains true if $\frac{2}{3}$ is replaced by any number less than 1.}
    
    \begin{problem}[1in]
    Suppose you have a series $\displaystyle \sum_{k=1}^\infty b_k$
    whose terms are all positive, and you know that $\displaystyle
    \lim_{k \to \infty} \frac{b_{k+1}}{b_k} = \frac{2}{3}$.  Does the
    series converge or diverge, or is there not enough information to
    tell? 
    \end{problem}

    \begin{solution}
    Since $\frac{b_{k+1}}{b_k}$ tends to $\frac{2}{3} = 0.666\ldots$, we can be sure that, eventually, $0.6 <
    \frac{b_{k+1}}{b_k} < 0.7$.  (Here, $0.6$ and $0.7$ are just
    arbitrary numbers smaller and bigger than $0.666\ldots$.)  Then, it
    is possible to compare directly to a geometric series with common
    ratio $0.7$.
    \end{solution}

    \item

    %\note{This was not on their homework, but it extends the previous problem to series with positive and negative terms.}
    
    \begin{problem}[2in]
    Suppose you have a series $\displaystyle \sum_{k=1}^\infty c_k$
    whose terms are nonzero (but the series could have both positive and
    negative terms), and you know that $\displaystyle \lim_{k \to
    \infty} \left| \frac{c_{k+1}}{c_k} \right| = \frac{2}{3}$.  Does
    the series $\displaystyle \sum_{k=1}^\infty |c_k|$ converge or
    diverge, or is there not enough information to tell?  What does this
    tell you about the series $\displaystyle \sum_{k=1}^\infty c_k$?
    \end{problem}

    \begin{solution}
    If we let $b_k = |c_k|$, then the series $\displaystyle
    \sum_{k=1}^\infty b_k$ has all positive terms, and we know that
    $\displaystyle \lim_{k \to \infty} \frac{b_{k+1}}{b_k} =
    \frac{2}{3}$.  This is exactly the situation described in
    {{\#1}(b)}, so we know that $\displaystyle
    \sum_{k=1}^\infty b_k$ converges.  Since $b_k = |c_k|$, we've shown
    that $\displaystyle \sum_{k=1}^\infty |c_k|$ converges.  This means
    that $\displaystyle \sum_{k=1}^\infty c_k$ is absolutely convergent
    (by definition of ``absolutely convergent'').  Every absolutely
    convergent series is convergent, so this also tells us that
    $\displaystyle \sum_{k=1}^\infty c_k$ converges. 
    \end{solution}
    
  \end{enumerate}

  \pspace{\newpage}

\begin{problemonly}
\begin{framed}
\textbf{Ratio Test:}  Suppose that $\displaystyle \sum a_{k}$ is a series whose terms are
nonzero, and such that $\displaystyle L = \lim_{k \to \infty} \left| \frac{a_{k+1}}{a_{k}} \right|$
either exists or is $\infty$.  Then
\begin{itemize}
  \item If $L < 1$, then the series converges absolutely.
  \item If $L > 1$ (which includes the case $L = \infty$) then the series diverges.
\end{itemize}
If $L = 1$, then the Ratio Test is inconclusive and you should try some other test.
\end{framed}
\end{problemonly}



  \item 
  What does the Ratio Test tell you about the following series?
  \begin{enumerate}
    \item
    
    \begin{problem}[0.9in]
    $\displaystyle \sum_{k=1}^\infty \frac{k}{3^k}$.
    \end{problem}

    \begin{solution}
    Let $a_k = \frac{k}{3^k}$.  Then,
    \begin{align*}
    \lim_{k \to \infty} \left| \frac{a_{k+1}}{a_k} \right| & =
    \lim_{k \to \infty} \left| \frac{(k + 1)}{3^{k+1}} \cdot \frac{3^k}{k}
    \right| \\
    & = \lim_{k \to \infty} \left| \frac{1}{3} \cdot \frac{k + 1}{k}
    \right| \\ 
    & = \frac{1}{3} \lim_{k \to \infty} \left| 1 + \frac{1}{k}
    \right| \\ 
    & = \frac{1}{3}
    \end{align*}
    Since this limit exists and is less than 1, the Ratio Test says that
    the series \fbox{converges absolutely}.
    \end{solution}
    
    \item

    \note{You might want to point out how well the Ratio Test works here
    with the powers and factorials.  This ``explains'' (later) why we so
    often use the Ratio Test when dealing with power series.}
    
    \begin{problem}[0.9in]
    $\displaystyle\sum_{k=0}^\infty (-1)^{k + 1} \frac{1000^k}{k!}$. 
    \end{problem}

    \begin{solution}
    Since $\displaystyle \lim_{k \to \infty} \left| \frac{(-1)^{(k + 1)
    + 1} \frac{1000^{k + 1}}{(k + 1)!}}{(-1)^{k + 1}
    \frac{1000^k}{k!}} \right| = \lim_{k \to \infty}
    \frac{1000^{k+1} \cdot k!}{(k + 1)! \cdot1000^k} = \lim_{k \to
    \infty} \frac{1000}{k + 1} = 0$, the Ratio Test says that the
    series \fbox{converges absolutely}.
    \end{solution}

    \item
    \begin{problem}[0.9in]
    $\displaystyle \sum_{k=1}^\infty \frac{1}{k}$.
    \end{problem}

    \begin{solution}
    $\displaystyle \lim_{k \to \infty} \left|
    \frac{\frac{1}{k+1}}{\frac{1}{k}} \right| = \displaystyle \lim_{k
    \to \infty} \left| \frac{k}{k + 1} \right| = \lim_{k \to \infty}
    \left| \frac{1}{1 + \frac{1}{k}} \right| = 1$, so the Ratio Test is
    \fbox{inconclusive}.  (Of course, we know the series diverges.)
    \end{solution}

    \item
    \begin{problem}[0.9in]
    $\displaystyle \sum_{k=1}^\infty (-1)^{k+1} \frac{1}{k}$.
    \end{problem}

    \begin{solution}
    $\displaystyle \lim_{k \to \infty} \left| \frac{(-1)^{k + 2}
    \frac{1}{k+1}}{(-1)^{k+1} \frac{1}{k}} \right| = \displaystyle
    \lim_{k \to \infty} \left| \frac{k}{k + 1} \right| = \lim_{k \to
    \infty} \left| \frac{1}{1 + \frac{1}{k}} \right| = 1$, so the
    Ratio Test is \fbox{inconclusive}.  (Of course, we know the series
    converges conditionally.)
    \end{solution}      
    
    \item
    \begin{problem}[0.9in]
    $\displaystyle \sum_{k=1}^\infty \frac{1}{k^2}$.
    \end{problem}

    \begin{solution}
    $\displaystyle \lim_{k \to \infty} \left|
    \frac{\frac{1}{(k+1)^2}}{\frac{1}{k^2}} \right| = \displaystyle
    \lim_{k \to \infty} \left| \frac{k^2}{(k + 1)^2} \right| = \lim_{k
    \to \infty} \left| \frac{1}{\left(1 + \frac{1}{k} \right)^2}
    \right| = 1$, so the Ratio Test is \fbox{inconclusive}.  (Of course,
    we know the series converges absolutely.)
    \end{solution}

  \end{enumerate}  
\pspace{\newpage}
    \item For each of the Taylor series below, use the Ratio Test to determine for which values of $x$ the series converges.
  \begin{enumerate}
    \item 
    \begin{problem}[\vfill]
    The Taylor series for $e^x$ centered at 0: $\displaystyle \sum_{k=0}^\infty \frac{x^k}{k!} = 1 + x + \frac{x^2}{2!} + \frac{x^3}{3!} + \cdots$
\end{problem}

    \begin{solution}
    The Ratio Test says we should look at
    \begin{align*}
    \lim_{k \to \infty} \left| \frac{\frac{x^{k+1}}{(k +
    1)!}}{\frac{x^k}{k!}} \right| 
    & = \lim_{k \to \infty} \left| \frac{x^{k+1}}{x^k} \cdot
    \frac{k!}{(k + 1)!} \right| \\
    & = \lim_{k \to \infty} \left| x \cdot \frac{1}{k + 1} \right|
    \\
    & = |x| \lim_{k \to \infty} \left| \frac{1}{k + 1} \right| \\
    & = |x| \cdot 0 \\
    & = 0
    \end{align*}
    Since the limit exists and is less than 1, the Ratio Test says that
    the series converges absolutely, no matter what $x$ is.  That is,
    the Taylor series \fbox{converges absolutely for all $x$}. 
    
    Note that we still don't know if the Taylor series for $e^x$ is actually \textit{equal} to $e^x$. All we know is that the Taylor series for $e^x$ converges to \textit{something} for all $x$ values.
    \end{solution}

\item
\begin{problem}[\vfill]
The Taylor series for $\sin(x)$ centered at 0: $\displaystyle \sum_{k=0}^\infty (-1)^k
  \frac{x^{2k+1}}{(2k+1)!}= x-\frac{1}{3!}x^3+\frac{1}{5!}x^5-\cdots$ 
  \end{problem}


\begin{solution}
To use the Ratio Test, we look at this limit: 
  \begin{align*}
   \lim_{k \to \infty} \left| \frac{(-1)^{k + 1}
   \frac{x^{2(k+1+1}}{(2(k +1)+ 1)!}}{(-1)^k \frac{x^{2k+1}}{(2k+1)!}}
   \right|  
   & = \lim_{k \to \infty} \left| \frac{x^{2k + 3} \cdot (2k+1)!}{(2k +
   3)! \cdot x^{2k+1}} \right| \\ 
   & = \lim_{k \to \infty} \left| \frac{x^2}{(2k + 3)(2k + 2)}
   \right| \\
   & = \left| x^2 \right| \lim_{k \to \infty} \left| \frac{1}{(2k +
   3)(2k + 2)} \right| \\
   & = \left| x^2 \right| \cdot 0 \\
   & = 0
   \end{align*}
   (Since we're taking the limit as $k$ tends to infinity, we treat $x$
   as a constant when taking the limit.)  Therefore, the Ratio Test says
   that, no matter what $x$ is, the series converges.  That is, the
   series \fbox{converges for all $x$}.
  
      Note that we still don't know if the Taylor series for $\sin(x)$ is actually \textit{equal} to $\sin(x)$. All we know is that the Taylor series for $\sin(x)$ converges to \textit{something} for all $x$ values.
\end{solution}
  \end{enumerate}

  \item 
  \note{Students often have horrendous difficulty with the ratio
  $\frac{a_{k+1}}{a_k}$ here.  Sometimes they will write down $a_{k+1}$
  incorrectly, and they also have trouble simplifying ratios of
  factorials.  If they have trouble, encourage them to try a specific
  value of $k$, like $k = 10$.}
  
  \begin{problem}[\vfill]
  What does the Ratio Test tell you about the series $\displaystyle
  \sum_{k=1}^\infty \frac{(3k)!}{5^{2k}}$?
  \end{problem}

  \begin{solution}
  Let $\displaystyle a_k = \frac{(3k)!}{5^{2k}}$.  The Ratio Test says
  that we should look at $\displaystyle \lim_{k \to \infty} \left|
  \frac{a_{k+1}}{a_k} \right|$:
  \begin{align*}
  \lim_{k \to \infty} \left| \frac{a_{k+1}}{a_k} \right| & = \lim_{k
  \to \infty} \left| \frac{\frac{[3(k + 1)]!}{5^{2 (k +
  1)}}}{\frac{(3k)!}{5^{2k}}} \right| \\
  & = \lim_{k \to \infty} \left|\frac{(3k + 3)!}{5^{2k + 2}}
  \cdot \frac{5^{2k}}{(3k)!} \right| \\
  & = \lim_{k \to \infty} \left| \frac{(3k + 3)(3k + 2)(3k +
  1)}{5^2} \right| \\
  & = \infty
  \end{align*}
  Since this limit is equal to $\infty$, the Ratio Test says that the
  series \fbox{diverges}.
  \end{solution}

\pspace{\newpage}

 \item \label{ln-interval-of-convergence-v2}  
   
Recall that the Taylor series for $\ln(1+x)$ centered at 0 is $\displaystyle
\sum_{k=1}^\infty (-1)^{k+1} \frac{x^k}{k}$. We aim to find for what values of $x$ does this series converge.

%We saw that this series converges when $x=1$ due to the Alternating Series Test for Convergence. However, we're not yet sure what happens for other values of $x$.
  
  \begin{enumerate}
    \item \label{ln-interval-of-convergence-v2-ratio-test}
    
    \begin{problem}[3in]
    What does the Ratio Test tell you about the convergence and divergence
    of this series?
    \end{problem}

    \begin{solution}
    \begin{align*}
    \lim_{k \to \infty} \left| \frac{(-1)^{(k + 1) +1}
    \frac{x^{k + 1}}{k + 1}}{(-1)^{k+1} \frac{x^k}{k}} \right| & 
    = \lim_{k \to \infty} \left| \frac{x^{k+1} \cdot k}{(k +
    1) \cdot x^k} \right| \\
    & = \lim_{k \to \infty} \left| x \cdot \frac{k}{k + 1} \right| \\
    & = |x| \lim_{k \to \infty} \left| \frac{k}{k + 1} \right| \\
    & = |x| \cdot 1 \\
    & = |x|
    \end{align*}
    So, \fbox{the Ratio Test says that the series converges when $|x| <
    1$ and diverges when $|x| > 1$}.  The Ratio Test is inconclusive
    when $|x| = 1$.
    \end{solution}

    \item \label{ln-interval-of-convergence-v2-endpoints}
    
\note{Students very often have trouble simplifying $(-1)^{k + 1} (-1)^k$,
    so keep an eye out for this.}

    
    \begin{problem}[2.5in]
    There are two values of $x$ for which the Ratio Test gives no
    information.  What are those two values?  Does the series converge for
    these two values of $x$? 
    \end{problem}

    \begin{solution} 
    As we said already, the Ratio Test doesn't tell us anything when
    $|x| = 1$, i.e., when \fbox{$x = \pm 1$}.

    If we plug $x = -1$ into the given Taylor series, we get
    $$\displaystyle \sum_{k=1}^\infty \frac{(-1)^{k + 1}}{k} (-1)^k =
    \sum_{k=1}^\infty \frac{(-1)^{2k + 1}}{k} = - \sum_{k=1}^\infty
    \frac{1}{k}\footnote{One way to think about this is that $2k + 1$
    is always an odd number, and $-1$ raised to an odd power is $-1$.
    Another way to think about it is to use exponentiation rules:
    $(-1)^{2k + 1} = (-1)^{2k} \cdot (-1)^1 = [(-1)^2]^k \cdot (-1) =
    1^k (-1) = -1$.  Or, you could write the series out in $\cdots$
    notation instead of sigma notation. The last way is perhaps the easiest way to see this fact.} $$ 

    This is just $-1$ times the
    harmonic series, and we know the harmonic series diverges.  So, the
    power series \fbox{diverges when $x = -1$}. 

    If we plug $x = 1$ into the given Taylor series, we get
    $\displaystyle \sum_{k=1}^\infty \frac{(-1)^{k + 1}}{k}$.  This is
    the alternating harmonic series, and we already showed that this series converges due to the Alternating Series Test. So, the Taylor
    series \fbox{converges when $x = 1$}. 
    \end{solution}

    \item
    \begin{problem}[2in]
    Putting \ref{ln-interval-of-convergence-v2-ratio-test} and
    \ref{ln-interval-of-convergence-v2-endpoints} together, what are
    \textit{all} values of $x$ for which the series $\displaystyle
    \sum_{k=1}^\infty (-1)^{k+1} \frac{x^k}{k}$ converges?
    \end{problem}

    \begin{solution}
    We've showed that the series converges when $-1 < x \leq 1$ (and
    diverges for all other values of $x$).
    
    This finally explains why on Weekly P-Set 5, the graphs of the Taylor polynomials for $\ln(1+x)$ centered at 0 didn't seem to exist outside of the domain $[-1,1]$, regardless of how high the degree of the polynomial was. Outside of the interval $(-1, 1]$, the Taylor series for $\ln(1+x)$ diverges, so the sums of the terms of the Taylor polynomials don't approach any value outside of this interval. Thus, for $x$-values outside of $(-1, 1]$, the graphs of the Taylor polynomials of $\ln(1+x)$ will never get closer to the function $\ln(1+x)$ as the degree of the polynomial increases.
    \end{solution}
\end{enumerate}




\end{enumerate}




\end{document}
